\subsection{Continuity and Regularity of Measures}

\begin{corollary} \label{cor:measure-continuity-up-down} \hlt{Measure Continuity}\\
Let $\mu$ be a measure. If $E_k\uparrow E$, then $\mu(E_k)\uparrow\mu(E)$. If $F_k\downarrow F$ and $\mu(F_1)<+\infty$, then $\mu(F_k)\downarrow\mu(F)$.
\begin{proof}
Disjointize $E_k\uparrow E$ by $E'_1=E_1$ and $E'_k=E_k\setminus E_{k-1}$; countable additivity makes $\mu(E_k)$ the $k$th partial sum for $\mu(E)$. For $F_k\downarrow F$, apply this to $F_1\setminus F_k\uparrow F_1\setminus F$ and subtract from finite $\mu(F_1)$.\\
\hlt{Structural Commentary}: Monotone set limits are disjoint series or complements of such series.\\
\hlt{Key Insight}: Downward continuity needs one finite anchor so subtraction is legal.
\end{proof}
\end{corollary}

\begin{remark} \label{rem:borel-null-regular-measurable-sets} \hlt{Borel Sets, Null Sets, and Regular Approximation}\\
Open, closed, and Borel subsets of $\R^p$ are measurable. Each $A\in\mathcal{M}(\mu)$ admits closed--open approximations $F\subseteq A\subseteq G$ with $\mu(G\setminus A)$ and $\mu(A\setminus F)$ arbitrarily small. In Lebesgue measure, one may take $F_\sigma\subseteq A\subseteq G_\delta$ with $\mu(G_\delta\setminus F_\sigma)=0$. Null sets form a $\sigma$-ring, and countable subsets of $\R^p$ are null.\\
\hlt{Structural Commentary}: Measurable sets are Borel sets modified by null sets.\\
\hlt{Key Insight}: Regularity explains why open and closed approximations are enough for integration.
\end{remark}

\begin{corollary} \label{cor:lebesgue-regularity-and-borel-completion-real-line} \hlt{Lebesgue Regularity and Borel Completion on $\R$}\\
For $A\in\mathcal{L}^1$,
\begin{enumerate}[label=(\roman*)]
\setlength{\itemsep}{0pt}
\item $m(A)=\inf\{m(G):A\subseteq G,\ G$ open in $\R\}$.
\item $m(A)=\sup\{m(K):K\subseteq A,\ K$ compact in $\R\}$, finite-measure case applied to $A\cap(-n,n]$; $m(A)=+\infty$.
\item There exists $B\in\mathcal{B}(\R)$ with $A\subseteq B$ and $m(B\setminus A)=0$; equivalently, $\mathcal{L}^1$ is the completion of $\mathcal{B}(\R)$.
\end{enumerate}
\begin{proof}
\begin{enumerate}[label=(\roman*)]
\setlength{\itemsep}{0pt}
\item For $\varepsilon>0$, regular approximation gives open $G\supseteq A$ with $m(G\setminus A)<\varepsilon$, hence $m(G)\leq m(A)+\varepsilon$.\\
Monotonicity gives the reverse inequality.
\item If $m(A)<+\infty$, choose closed $F\subseteq A$ with $m(A\setminus F)<\varepsilon$. Since $F\cap[-n,n]\uparrow F$, continuity from below gives compact $K=F\cap[-n,n]$ with $m(F)-m(K)<\varepsilon$, so $m(A)-m(K)<2\varepsilon$. If $m(A)=+\infty$, apply the finite case to $A_n=A\cap(-n,n]$ and let $n\to\infty$.
\item Choose open $G_k\supseteq A$ with $m(G_k\setminus A)<1/k$, and set $B=\bigcap_kG_k$. Then $B\in\mathcal{B}(\R)$, $A\subseteq B$, and $m(B\setminus A)=0$. Thus every Lebesgue set is Borel plus a null subset.\\
Conversely, Borel sets and null subsets are Lebesgue measurable.
\end{enumerate}
\hlt{Structural Commentary}: Lebesgue sets are Borel sets plus null-set ambiguity.\\
\hlt{Key Insight}: Open outer approximation and compact inner approximation are the usable form of regularity.
\end{proof}
\end{corollary}

\begin{proposition} \label{prop:real-line-local-measurability-unit-intervals} \hlt{Local Measurability on Unit Intervals}\\
Let $I_n\defeq(n,n+1]$, $n\in\Z$. For $E\subseteq\R$, $E\in\mathcal{L}^1$ iff $E\cap I_n\in\mathcal{L}^1$ for all $n\in\Z$.
\begin{proof}
\begin{enumerate}[label=\roman*.]
\setlength{\itemsep}{0pt}
\item[$(\Rightarrow)$] Since $I_n\in\mathcal{B}(\R)\subseteq\mathcal{L}^1$, $E\in\mathcal{L}^1$ implies $E\cap I_n\in\mathcal{L}^1$ for all $n$.
\item[$(\Leftarrow)$] If $E\cap I_n\in\mathcal{L}^1$ for all $n$, then $E=\bigcup_{n\in\Z}(E\cap I_n)\in\mathcal{L}^1$.
\end{enumerate}
\hlt{Structural Commentary}: A countable Borel exhaustion localizes measurability.\\
\hlt{Key Insight}: Local measurable pieces patch by countable unions.
\end{proof}
\end{proposition}

\begin{proposition} \label{prop:finite-measure-meets-every-unit-interval} \hlt{Finite Measure Meeting Every Unit Interval}\\
With $I_n\defeq(n,n+1]$, there exists $E\in\mathcal{B}(\R)$ such that $0<m(E\cap I_n)$ for all $n\in\Z$ and $m(E)<+\infty$.
\begin{proof}
Set $\varepsilon_n\defeq2^{-\abs{n}-3}$ and $E\defeq\bigcup_{n\in\Z}(n,n+\varepsilon_n]$. Then $(n,n+\varepsilon_n]\subseteq I_n$ are pairwise disjoint Borel intervals, so $E\in\mathcal{B}(\R)$ and $m(E\cap I_n)=\varepsilon_n>0$. Countable additivity gives $m(E)=\sum_{n\in\Z}\varepsilon_n=\frac38<+\infty$.\\
\hlt{Structural Commentary}: Finite measure can be spread over countably many cells.\\
\hlt{Key Insight}: Positive local masses can be summable globally.
\end{proof}
\end{proposition}

\begin{definition} \label{def:symmetric-difference} \hlt{Symmetric Difference}\\
For $A,B\subseteq X$, define $A\triangle B\defeq(A\setminus B)\cup(B\setminus A)=(A\cup B)\setminus(A\cap B)$.\\
Equivalently, $x\in A\triangle B$ iff $x$ belongs to exactly one of $A,B$.
\end{definition}

\begin{proposition} \label{prop:null-sets-sigma-ideal} \hlt{Null Sets Form a Sigma-Ideal}\\
Let $(X,\mathcal{M},\mu)$ be a measure space and $\mathcal{Z}\defeq\{N\in\mathcal{M}:\mu(N)=0\}$.\\
Then $\mathcal{Z}$ is a $\sigma$-ring, and $N\in\mathcal{Z}$, $E\in\mathcal{M}$ imply $N\cap E\in\mathcal{Z}$.
\begin{proof}
If $(N_k)\subseteq\mathcal{Z}$, then $\mu(\bigcup_kN_k)\leq\sum_k\mu(N_k)=0$, so $\bigcup_kN_k\in\mathcal{Z}$.\\
If $A,B\in\mathcal{Z}$, then $A\setminus B\subseteq A$, hence $A\setminus B\in\mathcal{Z}$. Finally, $N\cap E\subseteq N$ gives $N\cap E\in\mathcal{Z}$.\\
\hlt{Structural Commentary}: Null sets are stable under countable unions and measurable cuts.\\
\hlt{Key Insight}: Monotonicity makes every measurable piece of a null set null.
\end{proof}
\end{proposition}

\begin{proposition} \label{prop:completed-sigma-algebra} \hlt{Completed Sigma-Algebra}\\
Let $(X,\mathcal{M},\mu)$ be measure space, $\mathcal{Z}\defeq\{N\in\mathcal{M}:\mu(N)=0\}$; $\overline{\mathcal{M}}_\mu\defeq\{A\subseteq X:\exists E\in\mathcal{M},\ \exists N\in\mathcal{Z},\ A\triangle E\subseteq N\}$.
\begin{enumerate}[label=(\roman*)]
\setlength{\itemsep}{0pt}
\item $\mathcal{M}\subseteq\overline{\mathcal{M}}_\mu$, and $\overline{\mathcal{M}}_\mu$ is a $\sigma$-algebra;
\item $A\subseteq N\in\mathcal{Z}\Longrightarrow A\in\overline{\mathcal{M}}_\mu$;
\item $\overline{\mathcal{M}}_\mu=\{E\cup Z:E\in\mathcal{M},\ \exists N\in\mathcal{Z},\ Z\subseteq N\}$.
\end{enumerate}
\begin{proof}
\begin{enumerate}[label=(\roman*)]
\setlength{\itemsep}{0pt}
\item For $E\in\mathcal{M}$, $E\triangle E=\emptyset$, so $\mathcal{M}\subseteq\overline{\mathcal{M}}_\mu$. If $A\triangle E\subseteq N$, then $A^c\triangle E^c=A\triangle E\subseteq N$.\\
If $A_k\triangle E_k\subseteq N_k$, then $(\bigcup_kA_k)\triangle(\bigcup_kE_k)\subseteq\bigcup_kN_k\in\mathcal{Z}$. Hence $\overline{\mathcal{M}}_\mu$ is a $\sigma$-algebra.
\item If $A\subseteq N\in\mathcal{Z}$, then $A\triangle\emptyset\subseteq N$, so $A\in\overline{\mathcal{M}}_\mu$.
\item If $A\triangle E\subseteq N$, then $A=(E\setminus N)\cup(A\cap N)$, with $E\setminus N\in\mathcal{M}$ and $A\cap N\subseteq N$.\\
Conversely, if $A=E\cup Z$ and $Z\subseteq N\in\mathcal{Z}$, then $A\triangle E\subseteq Z\subseteq N$.
\end{enumerate}
\hlt{Structural Commentary}: Completion adds all null modifications of measurable sets.\\
\hlt{Key Insight}: Equality modulo null sets is the membership criterion.
\end{proof}
\end{proposition}

\begin{proposition} \label{prop:null-ideal-and-measure-completion} \hlt{Completed Measure}\\
Let $(X,\mathcal{M},\mu)$ be a measure space, $\mathcal{Z}\defeq\{N\in\mathcal{M}:\mu(N)=0\}$,\\
and $\overline{\mathcal{M}}_\mu\defeq\{A\subseteq X:\exists E\in\mathcal{M},\ \exists N\in\mathcal{Z},\ A\triangle E\subseteq N\}$. For $A\in\overline{\mathcal{M}}_\mu$, define $\overline{\mu}(A)\defeq\mu(E)$ whenever $E\in\mathcal{M}$ and $A\triangle E\subseteq N\in\mathcal{Z}$. Then $\overline{\mu}$ is a well-defined complete measure on $\overline{\mathcal{M}}_\mu$, and $\overline{\mu}|_{\mathcal{M}}=\mu$.
\begin{proof}
If $A\triangle E\subseteq N_1$ and $A\triangle F\subseteq N_2$, with $E,F\in\mathcal{M}$ and $N_1,N_2\in\mathcal{Z}$, then $E\triangle F\subseteq N_1\cup N_2$, so $\mu(E)=\mu(F)$. Hence $\overline{\mu}$ is well defined; since $E\triangle E=\emptyset$, $\overline{\mu}|_{\mathcal{M}}=\mu$.\\
For disjoint $(A_k)\subseteq\overline{\mathcal{M}}_\mu$, choose $E_k\in\mathcal{M}$ and $N_k\in\mathcal{Z}$ with $A_k\triangle E_k\subseteq N_k$, and set $N\defeq\bigcup_kN_k$.\\
Then $E_i\cap E_j\subseteq N$ for $i\neq j$. Put $F_1\defeq E_1$ and $F_k\defeq E_k\setminus\bigcup_{j<k}E_j$.\\
Then $(F_k)$ disjoint, $A_k\triangle F_k\subseteq N$, $(\bigcup_kA_k)\triangle(\bigcup_kF_k)\subseteq N$, so $\overline{\mu}\left(\bigcup_kA_k\right)=\mu\left(\bigcup_kF_k\right)=\sum_k\mu(F_k)=\sum_k\overline{\mu}(A_k)$.\\
If $A\subseteq B\in\overline{\mathcal{M}}_\mu$ and $\overline{\mu}(B)=0$, choose $E\in\mathcal{M}$ and $N\in\mathcal{Z}$ with $B\triangle E\subseteq N$.\\
Then $\mu(E)=0$, so $A\subseteq B\subseteq E\cup N\in\mathcal{Z}$; hence $A\in\overline{\mathcal{M}}_\mu$ and $\overline{\mu}(A)=0$.\\
\hlt{Structural Commentary}: The completed measure ignores null errors.\\
\hlt{Key Insight}: Null perturbations do not change countable additivity.
\end{proof}
\end{proposition}

\begin{definition} \label{def:set-liminf-limsup} \hlt{Limit Inferior and Superior of Sets}\\
For $(E_n)\subseteq\mathcal{P}(X)$, define $\limsup_nE_n\defeq\bigcap_{m=1}^{\infty}\bigcup_{n\geq m}E_n$ and $\liminf_nE_n\defeq\bigcup_{m=1}^{\infty}\bigcap_{n\geq m}E_n$.
\begin{enumerate}[label=(\roman*)]
\setlength{\itemsep}{0pt}
\item \hlt{Infinitely often}: $x\in\limsup_nE_n$ iff $\{n:x\in E_n\}$ is infinite;
\item \hlt{Eventually always}: $x\in\liminf_nE_n$ iff $\{n:x\notin E_n\}$ is finite;
\item \hlt{Set limit}: if $\limsup_nE_n=\liminf_nE_n$, write $\lim_nE_n$ for the common set.
\end{enumerate}
\end{definition}

\begin{proposition} \label{prop:set-limits-basic-laws} \hlt{Basic Laws for Set Limits}\\
Let $(E_n)\subseteq\mathcal{P}(X)$.
\begin{enumerate}[label=(\roman*)]
\setlength{\itemsep}{0pt}
\item $\emptyset\subseteq\liminf_nE_n\subseteq\limsup_nE_n\subseteq X$.
\item $E_n\uparrow E\Longrightarrow\liminf_nE_n=\limsup_nE_n=E=\bigcup_nE_n$.
\item $E_n\downarrow E\Longrightarrow\liminf_nE_n=\limsup_nE_n=E=\bigcap_nE_n$.
\item $(E_n)\subseteq\mathcal{M}\Longrightarrow\liminf_nE_n,\limsup_nE_n\in\mathcal{M}$.
\end{enumerate}
\begin{proof}
\begin{enumerate}[label=(\roman*)]
\setlength{\itemsep}{0pt}
\item Since $\bigcap_m\bigcup_{n\geq m}E_n\subseteq\bigcup_{n\geq1}E_n\subseteq X$ and $\bigcap_{n\geq m}E_n\subseteq\bigcup_{n\geq m}E_n$ for all $m$, the inclusions follow.
\item If $E_n\uparrow E$, then $\{n:x\in E_n\}$ is nonempty iff it is cofinite, hence $\liminf_nE_n=\limsup_nE_n=\bigcup_nE_n=E$.
\item If $E_n\downarrow E$, then $\{n:x\in E_n\}$ is infinite iff it is all of $\N$, hence $\liminf_nE_n=\limsup_nE_n=\bigcap_nE_n=E$.
\item If $(E_n)\subseteq\mathcal{M}$, then $\bigcup_{n\geq m}E_n,\bigcap_{n\geq m}E_n\in\mathcal{M}$ $\forall m$; countable closure gives $\limsup_nE_n,\liminf_nE_n\in\mathcal{M}$.
\end{enumerate}
\hlt{Structural Commentary}: Set limsup means infinitely often; set liminf means eventually always.\\
\hlt{Key Insight}: Translate the definitions into pointwise membership statements.
\end{proof}
\end{proposition}

\begin{proposition} \label{prop:measure-liminf-limsup-inequalities} \hlt{Measure Bounds for Set Limits}\\
Let $(E_n)$ be measurable in a measure space $(X,\mathcal{M},\mu)$.
\begin{enumerate}[label=(\roman*)]
\setlength{\itemsep}{0pt}
\item $\mu(\liminf_nE_n)\leq\liminf_n\mu(E_n)$.
\item $\mu(\bigcup_{n=1}^{\infty}E_n)<+\infty\Longrightarrow\limsup_n\mu(E_n)\leq\mu(\limsup_nE_n)$.
\end{enumerate}
\begin{proof}
\begin{enumerate}[label=(\roman*)]
\setlength{\itemsep}{0pt}
\item Set $A_m\defeq\bigcap_{n\geq m}E_n$. Then $A_m\uparrow\liminf_nE_n$ and $A_m\subseteq E_n$ for all $n\geq m$, so $\mu(A_m)\leq\inf_{n\geq m}\mu(E_n)$.\\
By continuity from below, $\mu(\liminf_nE_n)=\lim_m\mu(A_m)\leq\sup_m\inf_{n\geq m}\mu(E_n)=\liminf_n\mu(E_n)$.
\item Set $B_m\defeq\bigcup_{n\geq m}E_n$. $B_m\downarrow\limsup_nE_n$, $\mu(B_1)<+\infty$, $E_n\subseteq B_m$ $\forall n\geq m$, so $\sup_{n\geq m}\mu(E_n)\leq\mu(B_m)$.\\
By continuity from above, $\limsup_n\mu(E_n)=\inf_m\sup_{n\geq m}\mu(E_n)\leq\lim_m\mu(B_m)=\mu(\limsup_nE_n)$.
\end{enumerate}
\hlt{Structural Commentary}: Tail intersections increase to liminf; tail unions decrease to limsup.\\
\hlt{Key Insight}: Compare each tail envelope to the corresponding tail bounds of the measures.
\end{proof}
\end{proposition}
